% =========================================================================
% ISA Transactions -- BLINDED SUBMISSION
% Single file, 12 pt Times Roman, max 30 pages, double-anonymized.
% Author names, affiliations, acknowledgments and the identifying
% repository URL are removed from this file. They belong in the separate
% title page / cover letter.
% =========================================================================
\documentclass[12pt,a4paper]{article}

\usepackage[utf8]{inputenc}
\usepackage[T1]{fontenc}
\usepackage{mathptmx}          % Times Roman text and math
\usepackage[margin=2.5cm]{geometry}
\usepackage{setspace}
\onehalfspacing
\usepackage{lineno}
\linenumbers
\usepackage{graphicx}
\usepackage{booktabs}
\usepackage{amsmath,amssymb}
\usepackage{microtype}
\usepackage{caption}
\usepackage{float}
\usepackage{xcolor}
\usepackage{url}
\usepackage[hidelinks]{hyperref}

\graphicspath{{figures/}}
\captionsetup{font=small,labelfont=bf}

% =================================================================
% REVISION MARKERS -- REMOVE BEFORE SUBMISSION
% \TBD{...} marks a value that must be recomputed from the saved
% predictions before submission. Run scripts/postonset_analysis.py
% and replace every marker. Compile will show them in red.
% =================================================================
\newcommand{\TBD}[1]{\textcolor{red}{\textbf{[TBD:~#1]}}}

% Alternative titles that foreground the new contribution:
%  (a) "Fault-Onset Labeling Determines Reported Accuracy in the
%       Tennessee Eastman Benchmark: A Reproducible Study of Six
%       Classical Classifiers"
%  (b) "How Much of TEP Diagnostic Error Is a Labeling Artifact?
%       A Leakage-Controlled Benchmark of Classical Classifiers"


\title{\textbf{Fault-Onset Labeling Bounds Reported Accuracy in the Tennessee
Eastman Benchmark: A Leakage-Controlled Study of Six Classical Classifiers}}

\author{}
\date{}

\begin{document}
\maketitle

\begin{abstract}
Fault diagnosis in industrial processes is complicated by nonlinear variable
interactions, temporal dependence, and fault signatures that overlap normal
operation. This work presents a reproducible benchmark of six classical
machine-learning methods for multiclass diagnosis in the Tennessee Eastman
Process (TEP): logistic regression, decision tree, random forest,
histogram-based gradient boosting, linear support vector machine, and XGBoost.
The expanded TEP data comprised 15.33 million observations, 52 process
variables, normal operation, and 20 fault classes. To control leakage, data
were partitioned by complete simulation runs, and standardization parameters
were estimated using training data only. Hyperparameters were selected on a
validation partition, followed by repeated validation over five seeds. The
prespecified primary criterion was macro-F1. XGBoost ranked first in repeated
validation, with macro-F1 of $0.7683\pm0.0017$ and Matthews correlation
coefficient (MCC) of $0.7426$. After freezing the model and protocol, a single
evaluation was performed on the complete official test partition containing
10.08 million observations. Under the literal labels distributed with the
dataset, the final model achieved accuracy of 0.6754, macro-F1 of 0.7200, and
MCC of 0.6621. We further show that these figures are bounded by a structural
property of the benchmark itself. Because each fault is injected 160 samples
into every 960-sample official test run, literal labeling imposes a recall
ceiling of $800/960=0.8333$ on every fault class. Faults 6 and 7 attained
recalls of 0.8333 and 0.8332, that is, the ceiling itself. Re-evaluating only
the post-onset segment raised macro-F1 from 0.7200 to 0.8054, accuracy from 0.6754 to
0.7949, and MCC from 0.6621 to 0.7852, without refitting the model. Among the
1.6 million pre-onset test observations, 92.0\% were assigned to normal
operation or to the weakly observable faults 3, 9 and 15, which is the
physically correct response and is nonetheless scored as error under literal
labels. The effect is not specific to the selected model: seven classifiers
spanning single trees, ensembles and linear baselines all show a relative
accuracy gain between 0.142 and 0.150, against a structural bound of 0.1587. Rescoring the validation partition under the same protocol raised its macro-F1
from 0.7706 to 0.7953, so that the apparent validation-to-test generalization
gap of $+0.0483$ becomes $-0.0125$ under a symmetric comparison. The gap
therefore reflects the difference in mislabeled fraction between development
runs (4\%) and test runs (16.7\%) rather than a loss of generalization. Seventeen classes were diagnosed reliably, whereas normal
operation and faults 3, 9, and 15 remained mutually confounded, consistent with
their long-documented weak observability. Nonlinear tree ensembles provide
strong diagnosis for most TEP conditions, but aggregate metrics conceal both
class-specific limitations and labeling artifacts that are rarely reported.
\end{abstract}

\noindent\textbf{Keywords:} Tennessee Eastman Process; fault diagnosis;
machine learning; XGBoost; process monitoring; reproducible benchmark;
fault-onset labeling.

\section{Introduction}

Industrial processes operate through interacting physical, chemical, and
control mechanisms. A local disturbance can propagate through feedback loops,
alter several measured variables, and eventually compromise product quality,
equipment integrity, safety, or environmental performance. Fault detection and
diagnosis (FDD) therefore constitute central elements of modern process
monitoring \cite{venkatasubramanian2003review,yin2014review,ji2022review}.
Detection determines whether the process has departed from acceptable
operation, while diagnosis seeks to identify the underlying condition so that
operators or automated systems can respond appropriately.

Model-based diagnosis can be effective when reliable first-principles models
are available. However, plant-wide systems are frequently nonlinear,
high-dimensional, stochastic, and only partially observed. These properties
have motivated data-driven monitoring, in which statistical or machine-learning
models infer diagnostic boundaries from historical measurements
\cite{chiang2001book,yin2014review}. Classical multivariate methods such as
principal-component analysis (PCA), partial least squares (PLS), and Fisher
discriminant analysis established important foundations for process monitoring
\cite{chiang2000discriminant}. More recent studies increasingly employ tree
ensembles, support vector machines, neural networks, and sequence models to
represent nonlinear relationships and dynamic behavior.

The Tennessee Eastman Process (TEP), originally introduced by Downs and Vogel
\cite{downs1993tep}, remains one of the most widely used benchmarks in this
field. It represents a reactor--separator--recycle chemical process with 41
measured variables, 11 manipulated variables, and a diverse set of disturbances.
Its continued value arises from the coexistence of relatively distinctive
faults and difficult conditions whose trajectories overlap normal operation.
The expanded simulation data published by Rieth et al. \cite{rieth2017data}
provide hundreds of independent runs per condition and non-overlapping random
seeds for development and testing, enabling large-scale and reproducible
evaluation.

Despite extensive use of the TEP, reported results are often difficult to
compare. Studies may split temporally adjacent observations rather than complete
runs, tune models on the test set, report only accuracy, or omit class-level
behavior. Treating autocorrelated observations as independent experimental
replicates can produce optimistic estimates, while a high global score can hide
failures in specific operating conditions. These issues are especially relevant
when the normal state is confused with a fault, because false alarms can reduce
operator confidence and generate unnecessary interventions.

This study addresses these limitations through a leakage-controlled benchmark
of six classical supervised classifiers: logistic regression, decision tree,
random forest, histogram-based gradient boosting, linear support vector machine
(SVM), and XGBoost. The study asks five questions:

\begin{enumerate}
  \item Which classical model provides the strongest macro-averaged diagnostic
  performance under repeated validation?
  \item Does the selected model maintain its performance on a complete,
  independent official test partition?
  \item Which operating conditions account for the remaining diagnostic errors?
  \item How much of the reported error is attributable to the fault-onset
  labeling convention rather than to the classifier?
  \item Which model offers the best compromise between diagnostic performance
  and computational cost?
\end{enumerate}

The main contribution is an auditable protocol that separates model selection
from final testing, preserves complete simulation runs, quantifies variability
over five seeds, reports confidence intervals, and examines precision, recall,
F1, and confusion patterns for all 21 classes. Rather than presenting a single
headline score, the analysis identifies a compact ambiguity group that defines
the practical limit of the selected classifier.

A second contribution is methodological and, in our view, of wider consequence
for the TEP literature. The expanded dataset injects each fault after a fixed
warm-up interval, so a fraction of every faulty run is nominally normal yet
carries the fault label. This fraction differs between development runs (4\%)
and official test runs (16.7\%). We quantify how much of the reported error and
of the apparent generalization gap is produced by this convention alone, and we
report every headline metric under both the literal labeling and a post-onset
labeling. To our knowledge this decomposition is not routinely reported, which
may account for part of the dispersion among published TEP results.

\section{Related Work}

\subsection{Data-driven process monitoring}

Data-driven FDD spans latent-variable modeling, classification, clustering,
change detection, and temporal modeling. The review by Venkatasubramanian et
al. \cite{venkatasubramanian2003review} established a broad taxonomy covering
quantitative models, qualitative models, and process-history-based methods.
Chiang et al. \cite{chiang2000discriminant,chiang2001book} showed how PCA,
discriminant PLS, and Fisher discriminant analysis could support chemical-process
diagnosis. Later reviews emphasized that industrial data may be nonlinear,
dynamic, nonstationary, multimodal, and autocorrelated
\cite{yin2014review,ji2022review}. These properties motivate flexible nonlinear
classifiers, but also demand careful validation.

\subsection{Classical machine learning for TEP diagnosis}

Decision trees recursively partition the feature space into regions associated
with class labels \cite{breiman1984cart}. Random forests reduce the instability
of individual trees by combining bootstrap samples and randomized feature
subsets \cite{breiman2001rf}. Gradient boosting instead constructs an additive
ensemble in which successive learners correct residual errors
\cite{friedman2001gbm}. XGBoost extends this principle with regularized
objectives, efficient histogram-oriented computation, column subsampling, and
systems-level optimizations \cite{chen2016xgboost}. These characteristics make
tree ensembles attractive for TEP data, where variables interact nonlinearly.

Linear models provide important controls. Multinomial logistic regression tests
whether approximately linear decision surfaces are sufficient, whereas linear
SVMs maximize class separation through margin-based learning
\cite{cortes1995svm}. Their inclusion helps determine whether the added
complexity of nonlinear ensembles is justified. All implementations in this
study were based on scikit-learn \cite{pedregosa2011sklearn}, except XGBoost,
which used its reference implementation \cite{chen2016xgboost}.

\subsection{Evaluation beyond accuracy}

Accuracy can be misleading when performance varies substantially among classes.
Macro-F1 gives equal importance to every class by averaging class-specific F1
scores, while weighted-F1 weights them by support. Balanced accuracy averages
class recalls. MCC summarizes agreement between predictions and true labels and
remains informative for multiclass problems \cite{matthews1975mcc}. Confusion
matrices are indispensable because they reveal whether errors are diffuse or
concentrated in specific class pairs. The present study therefore uses
macro-F1 as the primary model-selection criterion and treats the other metrics
as complementary evidence.

\section{Materials and Methods}

\subsection{Tennessee Eastman Process data}

The TEP simulates a plant-wide reactor--separator--recycle system and includes
41 process measurements (XMEAS) and 11 manipulated variables (XMV), totaling 52
predictors \cite{downs1993tep}. The classification task contained 21 labels:
class 0 represented fault-free operation and classes 1--20 represented the
predefined fault conditions. The CSV distribution used in this work is a format
conversion of the expanded TEP simulation data of Rieth et al.
\cite{rieth2017data,kaggleTEP}.

The four source files contained 15,330,000 observations and 55 columns: the 52
process variables plus the identifiers \texttt{faultNumber},
\texttt{simulationRun}, and \texttt{sample}. Development runs contained 500
observations each. Official test runs contained 960 observations each. Every
class had 500 official test runs, yielding

\begin{equation}
21\times 500\times 960=10{,}080{,}000
\end{equation}

test observations. Each class therefore contributed 480,000 observations to
the final test.

\subsection{Fault-onset windows and labeling protocols}
\label{sec:onset}

The expanded dataset is generated by simulating each condition for 25 hours
(development) or 48 hours (official test) with a sampling interval of three
minutes, and by injecting the fault after a fixed warm-up interval: one hour in
faulty development runs and eight hours in faulty test runs
\cite{rieth2017data,rieth2018issues,reinartz2021extended}. In sample indices,
the fault is active from sample 21 in development runs and from sample 161 in
official test runs. Fault-free runs contain no injection.

Consequently, the labels distributed with the dataset are only conditionally
correct. In a faulty development run, samples 1--20 (4.0\% of the run) describe
nominal operation while carrying a fault label; in a faulty official test run,
samples 1--160 (16.7\% of the run) do so. Two consequences follow directly and
are, to our knowledge, seldom made explicit.

First, under an observation-level protocol with literal labels, the recall of
every fault class is bounded above by

\begin{equation}
R^{\max}_c=\frac{960-160}{960}=0.8333,\qquad c=1,\dots,20,
\label{eq:ceiling}
\end{equation}

whenever the classifier assigns pre-onset observations to some class other than
$c$. No such ceiling applies to class 0, whose test runs contain no injection.
Second, because the mislabeled fraction differs between development (4.0\%) and
test (16.7\%), a metric computed on the two partitions is not measuring the
same quantity, and part of any observed validation-to-test degradation is
attributable to this asymmetry rather than to generalization.

We therefore report all headline metrics under two labeling protocols:

\begin{itemize}
  \item \textbf{Literal labeling} (L): every observation retains the label
  distributed with the dataset. This is the protocol implicitly adopted by most
  published observation-level TEP results and is retained here for
  comparability.
  \item \textbf{Post-onset labeling} (P): pre-onset observations of faulty runs
  are excluded \emph{at evaluation time}, so that each scored observation
  carries a label consistent with the physical state of the plant. The model is
  not refitted: protocol P rescores the predictions of the same frozen model
  that was fitted under protocol L, including its 4\% of contaminated training
  observations. The comparison L versus P therefore isolates the labeling
  convention alone, and the reported post-onset figures are conservative, since
  a model additionally trained without the contaminated observations could only
  do better.
\end{itemize}

A third variant, in which pre-onset observations are excluded from fitting as
well, was not evaluated here; it would confound the labeling effect with a
change in the training distribution and is left for future work.

Protocol L is used for model selection, so that the frozen-model decision is
unaffected by the present analysis. Protocol P is reported alongside it in
Section~\ref{sec:onseteffect}. A third option, relabeling pre-onset
observations as class 0, was not adopted because it alters the class prior of
the normal condition and would confound the analysis of the ambiguity group.

\subsection{Leakage-controlled partitioning and preprocessing}

The original development runs were divided at the run level into training and
validation partitions using a fixed random seed and a stratified 80/20 split
within each source and class. This produced 8,400 training runs (400 per class),
2,100 validation runs (100 per class), and 10,500 official test runs (500 per
class). No run was assigned to more than one partition.

All 52 predictors were standardized using

\begin{equation}
z_{ij}=\frac{x_{ij}-\mu_j^{(\mathrm{train})}}
{\sigma_j^{(\mathrm{train})}},
\end{equation}

where $\mu_j^{(\mathrm{train})}$ and $\sigma_j^{(\mathrm{train})}$ were
estimated exclusively from the 4.2 million training observations. The fitted
transformation was then applied unchanged to validation and test data. This
procedure prevented information from the validation or test partitions from
influencing preprocessing.

\subsection{Candidate classifiers}

Six model families were compared. Hyperparameter tuning evaluated three
candidates per family, producing 18 configurations. The selected configurations
were subsequently frozen:

\begin{itemize}
  \item \textbf{Logistic regression}: $C=1.0$, maximum 1,000 iterations;
  \item \textbf{Decision tree}: unrestricted depth and minimum leaf size 1;
  \item \textbf{Random forest}: 400 trees, unrestricted depth, 50\% of features
  considered per split, and minimum leaf size 1;
  \item \textbf{HistGradientBoosting}: learning rate 0.05, 350 boosting
  iterations, and at most 31 leaf nodes;
  \item \textbf{Linear SVM}: stochastic-gradient hinge-loss classifier with
  $\alpha=10^{-5}$, maximum 2,000 iterations, and tolerance $10^{-3}$;
  \item \textbf{XGBoost}: 250 trees, maximum depth 6, learning rate 0.1,
  row subsampling 0.8, column subsampling 0.8, multiclass probabilistic
  objective, and histogram tree construction.
\end{itemize}

The tuning stage ranked configurations by macro-F1 and used MCC as a secondary
criterion. The official test partition was not accepted as an input by either
the tuning or repeated-validation programs.

The linear SVM was implemented as a stochastic-gradient surrogate rather than
as an exact margin solver, because quadratic-programming solvers were not
tractable at this sample size. Its score should therefore be read as a lower
bound on linear-margin performance, and the evidence for nonlinearity reported
in Section~\ref{sec:nonlinearity} rests primarily on the multinomial logistic
regression, which was optimized to convergence. Kernel SVMs were not evaluated
for the same computational reason.

\subsection{Repeated validation}

After selecting the best configuration for each model family, variability was
measured across seeds 2026--2030. For each seed, 40 complete runs per class were
sampled from training and independently from validation. Thus, every repetition
used 420,000 training and 420,000 validation observations while retaining the
full 500-observation trajectories of all selected runs.

For each metric, the mean, sample standard deviation, and 95\% confidence
interval were calculated across the five repetitions. With $n=5$, the interval
was

\begin{equation}
\bar{x}\ \pm\ t_{0.975,4}\frac{s}{\sqrt{5}}.
\end{equation}

Because only five seeds were used, these intervals quantify observed
seed-to-seed variability but should not be interpreted as definitive evidence
of pairwise statistical superiority. In particular, overlapping intervals for
the two leading models were interpreted conservatively.

In addition to the interval estimates, the six model families were compared
with the Friedman rank test over the five seeds treated as blocks, followed by
the Nemenyi post-hoc procedure \cite{demsar2006statistical}. Pairwise
comparisons used the Wilcoxon signed-rank test with Holm correction. The
critical difference of the Nemenyi procedure was computed as
$\mathit{CD}=q_{\alpha}\sqrt{k(k+1)/6N}$ with $k=6$ models and $N=5$ blocks.
Results are reported in Section~\ref{sec:selection}. With five blocks these
procedures have low power by construction, and they are reported to document
that the prespecified ranking rule was not substituted for an inferential
check, not to establish pairwise superiority.

Because a single frozen evaluation cannot express sampling uncertainty, the
official-test metrics were additionally accompanied by 95\% intervals obtained
by resampling complete runs with replacement (1{,}000 bootstrap replicates over
the 10{,}500 test runs). The run, not the observation, is the resampling unit,
consistent with the argument that temporally adjacent observations are not
independent replicates.

\subsection{Frozen final evaluation}

XGBoost candidate 1 was selected because it achieved the highest repeated-
validation macro-F1. The final seed was fixed at 2026 before test evaluation.
The final development set combined 40 complete runs per class from training and
40 from validation, totaling 840,000 observations. The frozen model was then
evaluated once on all 10,080,000 official test observations, without
subsampling or post-test adjustment.

Two earlier executions aborted during structural validation, before any model
fitting or metric computation, because of an incorrect run-length assertion;
only that assertion was changed, and the full incident log is available in the
repository.

The choice of 840,000 development observations, rather than the full 5,250,000
available, was imposed by the computational budget. To verify that this choice
does not materially affect the conclusions, a learning curve was computed on
the validation partition over $\{10,20,40,80,160\}$ complete runs per class:
\TBD{learning-curve result and saturation point}.

\subsection{Performance measures and reproducibility}
\label{sec:reproducibility}

The reported measures were accuracy, balanced accuracy, precision, recall,
class-specific F1, macro-F1, weighted-F1, MCC, confusion matrix, training time,
inference time, inference time per observation, and serialized model size.
Zero divisions in class-specific measures were assigned zero. Data and output
integrity were monitored with SHA-256 hashes and manifests. The final execution
used Python 3.12.13, scikit-learn 1.9.0, and XGBoost 3.4.0 on Linux.

\section{Results}

\subsection{Model selection under repeated validation}
\label{sec:selection}

Table~\ref{tab:validation} summarizes the five-seed validation results. XGBoost
ranked first with macro-F1 of $0.7683\pm0.0017$ and MCC of
$0.7426\pm0.0018$. HistGradientBoosting was a close second, with macro-F1 of
$0.7643\pm0.0034$. Random forest achieved 0.7398, followed by decision tree at
0.6714. The two linear baselines performed substantially worse: logistic
regression achieved 0.4685 and linear SVM achieved 0.4438.

\begin{table}[t]
\centering
\caption{Repeated-validation performance over five seeds. Macro-F1 was the
prespecified primary selection criterion.}
\label{tab:validation}
\small
\begin{tabular}{clccc}
\toprule
Rank & Model & Macro-F1 (mean $\pm$ SD) & 95\% CI & MCC (mean $\pm$ SD)\\
\midrule
1 & XGBoost & $0.7683\pm0.0017$ & [0.7661, 0.7705] & $0.7426\pm0.0018$\\
2 & HistGradientBoosting & $0.7643\pm0.0034$ & [0.7601, 0.7686] & $0.7391\pm0.0045$\\
3 & Random Forest & $0.7398\pm0.0021$ & [0.7372, 0.7424] & $0.7159\pm0.0021$\\
4 & Decision Tree & $0.6714\pm0.0013$ & [0.6699, 0.6730] & $0.6537\pm0.0015$\\
5 & Logistic Regression & $0.4685\pm0.0074$ & [0.4592, 0.4777] & $0.4567\pm0.0073$\\
6 & Linear SVM & $0.4438\pm0.0099$ & [0.4314, 0.4561] & $0.4377\pm0.0109$\\
\bottomrule
\end{tabular}
\end{table}

The Friedman test over the five seeds rejected the hypothesis of equal
performance among the six families ($\chi^2_{5}=24.54$, $p=1.71\times10^{-4}$),
with a Kendall coefficient of concordance of $W=0.982$. The ordering is thus
almost perfectly consistent from seed to seed: XGBoost obtained the best
macro-F1 in four of the five repetitions, and the mean ranks were 1.2
(XGBoost), 1.8 (HistGradientBoosting), 3.0 (random forest), 4.0 (decision
tree), 5.0 (logistic regression) and 6.0 (linear SVM).

Post-hoc resolution is nevertheless limited. With $N=5$ blocks the Nemenyi
critical difference is 3.372 mean ranks, so only three of the fifteen pairwise
comparisons are individually significant, all of them between a tree ensemble
and a linear baseline. The separation between XGBoost and
HistGradientBoosting is 0.6 mean ranks, far below the critical difference. The
Wilcoxon signed-rank test cannot help here either: with five paired
observations its smallest attainable two-sided $p$-value is 0.0625, so no
comparison can reach $\alpha=0.05$ even under a perfect win record, and after
Holm correction none does.

These two facts are complementary rather than contradictory. The high
concordance shows that the ranking is not an artifact of seed choice; the wide
critical difference shows that five repetitions cannot certify small
differences. Both support the conservative reading adopted here: the ordering
is reproducible, the gap between the two leading models is not resolvable at
this sample size, and additional seeds would be required to settle it.

Figure~\ref{fig:ranking} shows the ranking and its uncertainty. The confidence
intervals of XGBoost and HistGradientBoosting overlapped. Accordingly, XGBoost
was selected by the prespecified macro-F1 rule, but the result is not presented
as proof of a statistically significant difference between the two models.
The large separation between nonlinear ensembles and linear baselines indicates
that nonlinear interactions are central to this diagnostic task.

\begin{figure}[t]
\centering
\includegraphics[width=0.82\textwidth]{validation_model_ranking.pdf}
\caption{Repeated-validation macro-F1 for the six candidate algorithms. Error
bars denote 95\% confidence intervals across five seeds.}
\label{fig:ranking}
\end{figure}

\subsection{Official test performance}

After the protocol was frozen, the selected XGBoost model was fitted to 840,000
development observations and evaluated on the complete official test set.
Table~\ref{tab:test} presents the final metrics under literal labeling.
Accuracy and balanced accuracy were both 0.6754, macro-F1 and weighted-F1 were
both 0.7200, and MCC was 0.6621. The equality of accuracy and balanced
accuracy, and of macro-F1 and weighted-F1, follows from the exactly balanced
class supports, and holds only under literal labeling; under the post-onset
protocol of Section~\ref{sec:onseteffect} the fault classes retain 400,000
observations while the fault-free class retains 480,000, so the two pairs of
metrics separate. Macro-F1 exceeds accuracy because accuracy equals macro-recall
under balanced supports, while precision exceeds recall in most classes; the
harmonic mean of a high precision and a moderate recall is larger than the
recall alone. The gap between macro-F1 and accuracy is therefore itself a
signature of the systematic under-recall documented in
Section~\ref{sec:onseteffect}.

\begin{table}[H]
\centering
\caption{Frozen XGBoost performance on the official test set.}
\label{tab:test}
\small
\begin{tabular}{lc}
\toprule
Metric & Value\\
\midrule
Accuracy & 0.6754\\
Balanced accuracy & 0.6754\\
Macro-F1 & 0.7200\\
Weighted-F1 & 0.7200\\
MCC & 0.6621\\
Training time (s) & 170.34\\
Inference time (s) & 56.57\\
Inference (ms/observation) & 0.005612\\
Development observations & 840,000\\
Test observations & 10,080,000\\
\bottomrule
\end{tabular}
\end{table}

Relative to repeated validation, macro-F1 decreased by 0.0483 and MCC by
0.0805. Three mechanisms could account for this difference: independent random
states in the test simulations, the longer duration of test runs, and the
larger fraction of mislabeled pre-onset observations. The third mechanism is
both the largest and the only one that is purely a labeling artifact. Under
literal labeling the attainable recall per fault class falls from $480/500=0.96$
in development runs to $800/960=0.8333$ in test runs, a ratio of 0.868, which
is of the same order as the observed degradation. Section~\ref{sec:onseteffect}
shows that the gap does not merely shrink under post-onset labeling but changes
sign: the post-onset test macro-F1 of 0.8054 exceeds the literal-label
validation figure of 0.7683 by 0.0371.

That comparison alone would not be symmetric, because the validation figure
was itself computed under literal labels and carries its own 4\% pre-onset
contamination. The validation partition was therefore rescored under both
protocols as well. Applying the selected XGBoost configuration to all
1,050,000 validation observations gave a literal-label macro-F1 of 0.7706,
which agrees with the five-seed repeated-validation figure of 0.7683 to within
0.0023 and confirms that this partition is a faithful reference. Excluding the
40,000 pre-onset observations, 3.81\% of the partition, raised macro-F1 to
0.7953, accuracy from 0.7556 to 0.7843, and MCC from 0.7440 to 0.7739.

The correction on the validation partition is thus $+0.0246$ in macro-F1,
against $+0.0854$ on the test partition, a ratio of 0.29 that tracks the ratio
of contaminated fractions (3.81\% against 15.87\%). Transferring the measured
validation correction to the repeated-validation reference gives a post-onset
validation macro-F1 of approximately 0.7929, against a post-onset test macro-F1
of 0.8054. Table~\ref{tab:gap} summarizes the decomposition.

\begin{table}[H]
\centering
\caption{Validation-to-test difference in macro-F1 under both labeling
protocols. The gap present under literal labeling does not survive a symmetric
comparison.}
\label{tab:gap}
\small
\begin{tabular}{lccc}
\toprule
Protocol & Validation & Test & Difference\\
\midrule
Literal (L)    & 0.7683 & 0.7200 & $+0.0483$\\
Post-onset (P) & 0.7929 & 0.8054 & $-0.0125$\\
\bottomrule
\end{tabular}
\end{table}

Under a symmetric comparison the gap is not merely reduced but reversed, and
its magnitude falls to roughly a quarter of the literal-label value, within the
range of the seed-to-seed variability of Table~\ref{tab:validation}. We
therefore conclude that the reported generalization gap is an artifact of the
labeling convention rather than evidence of degradation on unseen trajectories,
and that the literal-label figure of 0.7200 understates the discriminative
capacity of the frozen model by 0.0854 in macro-F1.

Two caveats attach to this conclusion. The validation rescoring used the tuned
configuration fitted to the full training partition, not the five per-seed
refits of the repeated-validation procedure, so the transferred correction is
an estimate of the effect rather than an exact recomputation; the close
agreement between 0.7706 and 0.7683 bounds the error of that substitution.
Furthermore, the pre-onset destinations differ between partitions: 61.1\% of
pre-onset validation observations were assigned to the group $\{0,3,9,15\}$
against 92.0\% in the test partition, which is consistent with the development
segment being only 20 samples long and overlapping the simulation transient.

\subsection{Class-specific performance}

Figure~\ref{fig:classmetrics} and Table~\ref{tab:classwise} demonstrate marked
heterogeneity. Seventeen classes achieved F1 scores from 0.7429 to 0.9080.
Classes 6 and 7 were strongest, with F1 of 0.9080 and 0.9075, respectively.
Classes 1, 2, 5, and 14 also approached or exceeded 0.8958.

\begin{figure}[t]
\centering
\includegraphics[width=0.96\textwidth]{official_test_class_metrics.pdf}
\caption{Class-wise precision, recall, and F1 on the complete official test
partition. The dashed line indicates overall macro-F1. Class 0 is normal
operation.}
\label{fig:classmetrics}
\end{figure}

Four classes were persistently difficult: class 9 (F1=0.1635), class 15
(0.1683), normal operation (class 0; 0.1718), and class 3 (0.2357). Class 3 had
the highest recall within this group (0.4218) but very low precision (0.1636),
indicating systematic overprediction.

A pattern in the recall column deserves emphasis. Classes 6 and 7 reached
recalls of 0.8333 and 0.8332, which coincide with the structural ceiling of
Eq.~\eqref{eq:ceiling} to within one part in ten thousand. Classes 1, 2, 4, 5
and 14 lie between 0.8187 and 0.8243, that is, at 98--99\% of that ceiling. The
frozen model is therefore essentially saturated on these classes once the
pre-onset segment is removed, and their reported F1 values of roughly 0.88--0.91
are limited by the labeling convention rather than by the classifier.

\begin{table}[t]
\centering
\caption{Class-wise results for the frozen XGBoost model on the official test
set. Every class had support of 480,000 observations.}
\label{tab:classwise}
\small
\begin{tabular}{crrr@{\qquad}crrr}
\toprule
Class & Precision & Recall & F1 & Class & Precision & Recall & F1\\
\midrule
0  & 0.1481 & 0.2045 & 0.1718 & 11 & 0.8447 & 0.7095 & 0.7712\\
1  & 0.9894 & 0.8243 & 0.8993 & 12 & 0.9271 & 0.7669 & 0.8395\\
2  & 0.9896 & 0.8209 & 0.8974 & 13 & 0.9828 & 0.7444 & 0.8471\\
3  & 0.1636 & 0.4218 & 0.2357 & 14 & 0.9889 & 0.8187 & 0.8958\\
4  & 0.9477 & 0.8222 & 0.8805 & 15 & 0.1465 & 0.1978 & 0.1683\\
5  & 0.9879 & 0.8234 & 0.8982 & 16 & 0.9756 & 0.6681 & 0.7931\\
6  & 0.9973 & 0.8333 & 0.9080 & 17 & 0.9624 & 0.7809 & 0.8622\\
7  & 0.9964 & 0.8332 & 0.9075 & 18 & 0.9367 & 0.7760 & 0.8488\\
8  & 0.9665 & 0.7779 & 0.8620 & 19 & 0.7701 & 0.7559 & 0.7630\\
9  & 0.1328 & 0.2126 & 0.1635 & 20 & 0.8391 & 0.7029 & 0.7650\\
10 & 0.8066 & 0.6885 & 0.7429 & & & &\\
\bottomrule
\end{tabular}
\end{table}

\subsection{Effect of the fault-onset window}
\label{sec:onseteffect}

Table~\ref{tab:onset} compares the frozen model under the two labeling
protocols of Section~\ref{sec:onset}. No refitting, retuning or reselection was
performed; the same frozen model and the same predictions are scored under two
label conventions.

\begin{table}[t]
\centering
\caption{Frozen XGBoost under literal (L) and post-onset (P) labeling on the
official test partition. Intervals are 95\% run-level bootstrap intervals over
1,000 replicates.}
\label{tab:onset}
\small
\begin{tabular}{lcc}
\toprule
Metric & Literal (L) & Post-onset (P)\\
\midrule
Accuracy          & 0.6754 [0.6716, 0.6797] & 0.7949 [0.7898, 0.8006]\\
Balanced accuracy & 0.6754 [0.6748, 0.6760] & 0.8006 [0.7998, 0.8013]\\
Macro-F1          & 0.7200 [0.7192, 0.7208] & 0.8054 [0.8044, 0.8063]\\
MCC               & 0.6621 [0.6579, 0.6666] & 0.7852 [0.7798, 0.7910]\\
Observations      & 10,080,000 & 8,480,000\\
\bottomrule
\end{tabular}
\end{table}

Excluding the 1,600,000 pre-onset observations, 15.87\% of the partition,
raised macro-F1 by 0.0854, accuracy by 0.1195 and MCC by 0.1231. No model was
refitted and no hyperparameter was changed.

Two diagnostics establish that this is a labeling effect rather than a
favourable choice of subset. First, the frozen model assigned 16.10\% of the
pre-onset observations to class 0 and 75.86\% to classes 3, 9 or 15, that is,
92.0\% to the group of conditions that are indistinguishable from nominal
operation, and only 4.20\% to the fault label they nominally carry. The model
therefore identifies the pre-onset segment as a period in which nothing
detectable is occurring, which is correct, and literal labeling scores this as
error. Second, restricting the evaluation to post-onset observations raised the
recall of the 17 well-separated classes to between 0.8015 and 0.9999, the upper
value corresponding to essentially complete recovery of faults 6 and 7, while
the ambiguity group improved only marginally (F1 of 0.2217, 0.3238, 0.2100 and
0.2106 for classes 0, 3, 9 and 15, against 0.1718, 0.2357, 0.1635 and 0.1683
under literal labels). The ambiguity group is therefore not a labeling
artifact; the ceiling on the remaining seventeen classes is.

The two phenomena are nevertheless connected. Class 3 was predicted 1,237,562
times against a true support of 480,000, an excess of 757,562 observations. A
substantial part of that excess consists of pre-onset segments belonging to
other faults, which the model routes to the weakly observable conditions
because, at those instants, nothing has yet happened. The overprediction of
class 3 is thus partly a consequence of the labeling convention and partly a
consequence of genuine overlap, and only the post-onset analysis separates the
two.

This distinction matters for comparability. Published observation-level results
on the expanded TEP dataset do not consistently state whether pre-onset samples
were retained, discarded, or relabeled, yet the three choices are separated by
several points of macro-F1 for an identical model. We therefore recommend that
the treatment of the onset window be reported explicitly as part of the
experimental protocol.

\subsection{Generality of the labeling effect across model families}
\label{sec:generality}

If the effect of Section~\ref{sec:onseteffect} is a property of the labeling
convention rather than of the selected classifier, it should appear in every
model evaluated on the same partition. To test this, an auxiliary experiment
was rescored under both protocols. That experiment predates the frozen
benchmark, uses a per-class subsample of the test partition, and includes an
extremely randomized trees model in place of XGBoost; its absolute levels are
therefore not comparable with Table~\ref{tab:test}, and only the differences
between protocols are interpreted. No model was refitted.

\begin{table}[t]
\centering
\caption{Effect of post-onset labeling across model families. The last column
reports the relative accuracy gain $\Delta A/A_P$, whose structural upper bound
is $1-8{,}480{,}000/10{,}080{,}000=0.1587$.}
\label{tab:generality}
\small
\begin{tabular}{lccc}
\toprule
Model & Macro-F1 (L) & Macro-F1 (P) & $\Delta A/A_P$\\
\midrule
Decision tree           & 0.6330 & 0.7167 & 0.148\\
Extra trees             & 0.6965 & 0.7786 & 0.148\\
Random forest           & 0.7081 & 0.7912 & 0.149\\
HistGradientBoosting    & 0.7128 & 0.7972 & 0.149\\
XGBoost$^{\dagger}$     & 0.7200 & 0.8054 & 0.150\\
Logistic regression     & 0.4283 & 0.4698 & 0.142\\
Linear SVM              & 0.4207 & 0.4655 & 0.142\\
\bottomrule
\end{tabular}

\vspace{2pt}
{\footnotesize $^{\dagger}$Frozen model of Table~\ref{tab:test}, complete test
partition.}
\end{table}

The five tree-based models gained between 0.0821 and 0.0854 in macro-F1, a
spread of 0.003 across estimators whose capacities differ by orders of
magnitude, from a single unpruned tree to a regularized boosting ensemble. No
overfitting mechanism produces that uniformity, because the degradation
associated with overfitting scales with model capacity.

The relative accuracy gain explains the pattern quantitatively. If a classifier
assigns pre-onset observations to any class other than the labeled fault, its
literal accuracy is its post-onset accuracy scaled by the retained fraction of
the partition, so that $\Delta A/A_P$ is bounded above by 0.1587 independently
of the model. The seven models fall between 0.142 and 0.150, the residual being
the small fraction of pre-onset observations that match the literal label by
chance. The linear baselines gain less in absolute terms, 0.0415 and 0.0448,
simply because they start lower; in relative terms they are affected almost
identically. They also route far fewer pre-onset observations to the group
$\{0,3,9,15\}$, 20.6\% and 41.9\% against 65.7--70.2\% for the ensembles, and
are penalized in the same proportion regardless.

The consequence is that the effect is not a property of any estimator. Any
observation-level result reported on this dataset under literal labels
understates the post-onset discriminative capacity of the model by a factor
predictable from its own accuracy. Comparisons between published results are
correspondingly affected whenever the treatment of the onset window differs or
is unstated.

\subsection{Structure of the remaining errors}

The normalized confusion matrix in Fig.~\ref{fig:confusion} shows that errors
were not uniformly distributed. They concentrated in the four-class group
$\{0,3,9,15\}$. Among true class-9 observations, 32.39\% were predicted as
class 3, 20.24\% as class 0, and 18.08\% as class 15. Normal observations were
assigned to classes 3, 9, and 15 in 32.02\%, 21.51\%, and 18.24\% of cases,
respectively. For class 15, the corresponding proportions were 30.96\%,
20.53\%, and 20.03\% for predicted classes 3, 9, and 0.

\begin{figure}[H]
\centering
\includegraphics[width=\columnwidth]{official_test_confusion_matrix_normalized.pdf}
\caption{Row-normalized confusion matrix for the frozen XGBoost model on the
official test partition.}
\label{fig:confusion}
\end{figure}

Class 3 was predicted 1,237,562 times despite having 480,000 true observations.
This excess explains its low precision and demonstrates that the model used
class 3 as a frequent destination for ambiguous trajectories. The persistence
of the same difficult group during repeated validation and official testing
suggests an intrinsic overlap under an observation-level representation rather
than an isolated sampling artifact.

\subsection{Computational performance}

Table~\ref{tab:cost} reports the cost of all six families under identical
conditions during repeated validation, which is required to answer the
performance--cost question and cannot be inferred from the selected model
alone. Training times are reported as means over the five seeds and are
not directly comparable with the 170.34\,s of Table~\ref{tab:test}, which was
measured during the final evaluation on twice as many observations. The
degree of parallelism was an invocation-time argument of each stage rather
than a fixed setting, and it was not recorded in the run manifests, so the
absolute times of the two stages reflect different thread counts on a 24-core
host. The comparison between families within Table~\ref{tab:cost} is
unaffected, because those six runs shared a single invocation, but the
cross-stage comparison is not meaningful and no conclusion is drawn from it.
Recording the degree of parallelism alongside the timing is a reproducibility
requirement that this benchmark did not initially meet, and it has been added
to the released code.

\begin{table}[t]
\centering
\caption{Performance and computational cost of the six families under repeated
validation (mean over five seeds, 420,000 training observations per
repetition). Times measured on the environment of
Section~\ref{sec:reproducibility}.}
\label{tab:cost}
\small
\begin{tabular}{lrrrrr}
\toprule
Model & Macro-F1 & MCC & Train (s) & Inference ($\mu$s/obs) & Size (MB)\\
\midrule
XGBoost              & 0.7683 & 0.7426 & 864.7 & 7.03 & 18.9\\
HistGradientBoosting & 0.7643 & 0.7391 & 73.5 & 29.40 & 16.9\\
Random Forest        & 0.7398 & 0.7159 & 477.3 & 20.29 & 8703.8\\
Decision Tree        & 0.6714 & 0.6537 & 54.0 & 0.24 & 28.8\\
Logistic Regression  & 0.4685 & 0.4567 & 65.4 & 0.15 & 0.005\\
Linear SVM           & 0.4438 & 0.4377 & 4.1 & 0.16 & 0.006\\
\bottomrule
\end{tabular}
\end{table}

Final XGBoost training required 170.34 seconds. Prediction over 10.08 million
observations required 56.57 seconds, corresponding to 0.005612 ms per
observation or approximately 178,000 observations per second. The serialized
final model occupied approximately 6.8 MB. These values indicate that the
selected classifier is computationally feasible for high-throughput offline
analysis and potentially for online scoring, although end-to-end deployment
latency must also include data acquisition, preprocessing, communication, and
decision logic.

\section{Discussion}

\subsection{Nonlinearity and model ranking}
\label{sec:nonlinearity}

The ranking provides strong empirical evidence that the TEP classification
boundary is not adequately represented by the evaluated linear models.
XGBoost, HistGradientBoosting, and random forest substantially outperformed
logistic regression and linear SVM even though all models received the same
standardized predictors and run-level partitions. Tree ensembles can represent
threshold effects and high-order interactions without requiring an explicit
feature expansion, which is consistent with the nonlinear and feedback-driven
nature of the TEP.

The difference between XGBoost and HistGradientBoosting was small: 0.0040 in
mean macro-F1, with overlapping 95\% confidence intervals. Thus, the defensible
conclusion is that XGBoost won under the prespecified ranking rule, not that it
was conclusively superior under all resampling regimes. Table~\ref{tab:cost}
makes the practical consequence explicit, and it does not favour a single
model. HistGradientBoosting trained 11.8 times faster than XGBoost (73.5\,s
against 864.7\,s) for a deficit of 0.0040 in macro-F1 with overlapping
intervals, so it dominates whenever models are retrained frequently or the
training budget is constrained. The ordering reverses at inference time, where
XGBoost required 7.03\,$\mu$s per observation against 29.40\,$\mu$s, a factor of
4.2 in favour of XGBoost. Neither model dominates the other: the choice follows
from whether the deployment is retraining-bound or scoring-bound, and both
conclusions are invisible if only the selected model is profiled.

Model size adds a third and largely independent constraint. The random forest
required approximately 8.7\,GB, some 460 times the XGBoost artifact, for
0.0285 less macro-F1. Whatever its statistical merits, it is not a candidate
for embedded or edge deployment, and this would not be apparent from predictive
metrics alone. The two linear baselines occupied roughly 5\,kB and inferred in
0.15\,$\mu$s per observation, which quantifies precisely what is being paid for
the nonlinear gain of about 0.30 in macro-F1.

\subsection{Why global performance is insufficient}

The official macro-F1 of 0.7200 summarizes useful performance, but it does not
describe a uniformly reliable 21-class diagnostic system. If only macro-F1 or
accuracy had been reported, the strong results for 17 classes would obscure
the near-failure of normal operation and faults 3, 9, and 15. In operational
terms, low precision for normal operation and extensive exchange between normal
and faulty labels can increase false alarms and missed interventions.

This result illustrates why class-level precision, recall, F1, and confusion
structure should accompany aggregate measures in industrial diagnosis. It also
supports a distinction between \emph{benchmark discrimination} and
\emph{deployment readiness}. The present model is a strong benchmark
classifier for most faults, but additional safeguards are necessary before it
could support autonomous decisions.

\subsection{Interpretation of the ambiguity group}

The identity of the difficult group is not a new finding, and it would be
misleading to present it as one. Faults 3, 9 and 15 have been reported as
weakly detectable since the earliest systematic studies of the TEP: monitoring
statistics based on PCA, discriminant PLS and Fisher discriminant analysis
report missed-detection rates close to those of nominal operation for exactly
these conditions \cite{chiang2000discriminant,chiang2001book}, and later
comparative studies reproduce the pattern across a wide range of data-driven
statistics \cite{yin2012comparison}. The contribution here is the scale and the
protocol: the same group is recovered under run-level partitioning, on 10.08
million independent test observations, and it persists after the labeling
artifact of Section~\ref{sec:onseteffect} is removed. This strengthens the
interpretation that the overlap is a property of the measured variables rather
than of any particular estimator.

The mechanism, however, constrains which remedies can work. If these
disturbances leave no persistent signature in the 52 measured variables, then
temporal aggregation over a run reduces the variance of an estimate that is
already centered on an ambiguous region: it will sharpen the decision without
moving it toward the correct class. Run-level aggregation should therefore be
expected to help the 17 separable classes, whose pre-onset transients are
transient, far more than the group $\{0,3,9,15\}$, and this asymmetry is itself
a testable prediction. A hierarchical design that isolates the ambiguous group
and applies a discriminator with explicit memory to classes 0, 3, 9 and 15 is
the more promising direction, but its plausibility rests on whether any
signature exists at all; establishing this is a question of observability, not
of classifier capacity, and it is better addressed with residual-based or
first-principles analysis than with additional supervised tuning.

Both proposals arise from test-set diagnosis and must therefore be evaluated
under a newly defined development/test protocol; the present official test must
not be reused for tuning them.

\subsection{Threats to validity}

Several limitations delimit the conclusions. First, TEP is a high-fidelity
simulation benchmark, not a physical plant. Results may not transfer directly
to sensor drift, maintenance changes, unmodeled disturbances, communication
losses, or evolving operating modes in real facilities. Second, the benchmark
evaluated pointwise multiclass diagnosis and did not measure detection delay,
time-to-alarm, false alarms per run, or sustained alarm logic. Third, five seeds
provide a useful but small basis for uncertainty estimation. The resulting
$t$ intervals characterize the implemented experiment rather than all possible
data partitions.

Fourth, hyperparameter search was intentionally limited to three candidates per
model family to preserve computational feasibility and transparency. A larger
search might improve individual algorithms, particularly the linear baselines
or random forest. Fifth, the final model was fitted to 840,000 of the 5,250,000
available development observations, and model selection was carried out at
420,000; a ranking established at one training scale need not be invariant to
scale, and the learning curve reported in Section~\ref{sec:selection} bounds but
does not eliminate this concern. Sixth, the linear SVM was a stochastic-gradient
surrogate and its result should not be read as the best attainable
linear-margin performance. Seventh, the degree of parallelism was not recorded in the run manifests, so
the training times of Table~\ref{tab:cost} support ordinal comparison between
families, which shared one invocation, but not comparison against the timings
of other stages. Eighth, interpretation focused on output behavior
and did not include feature attribution or causal process analysis. Finally, the
dataset had balanced class support, whereas real industrial faults are rare and
highly imbalanced. Metrics, thresholds, and alarm costs would require
recalibration under field prevalence.

\section{Conclusion}

This work presented a reproducible classical-machine-learning benchmark for
21-class fault diagnosis in the Tennessee Eastman Process. Run-level splitting,
train-only standardization, separate tuning and validation, five-seed repetition,
and a single frozen official test evaluation were used to reduce leakage and
post-selection bias.

XGBoost achieved the highest repeated-validation macro-F1
($0.7683\pm0.0017$) and produced an official test macro-F1 of 0.7200, accuracy
of 0.6754, and MCC of 0.6621 over 10.08 million observations under literal
labeling, rising to a macro-F1 of 0.8054, an accuracy of 0.7949 and an MCC of
0.7852 once the pre-onset segment of each faulty run is excluded, with no
refitting of any kind. The validation partition, rescored under the same
protocol, moved from 0.7706 to 0.7953, so that the apparent generalization gap
of $+0.0483$ in macro-F1 becomes $-0.0125$ once the comparison is made
symmetric. Seventeen classes were diagnosed with F1 between
0.7429 and 0.9080, and several of them were shown to sit at the structural
recall ceiling imposed by the labeling convention rather than at the limit of
the classifier. Normal operation and faults 3, 9, and 15 remained strongly
confounded, in agreement with two decades of TEP monitoring literature, and
this confounding survives the removal of the labeling artifact.

Three conclusions follow. Nonlinear tree ensembles provide effective diagnosis
for most TEP conditions. Aggregate scores overstate practical reliability when
a small group of classes dominates the remaining errors. And a benchmark result
on the expanded TEP dataset is not interpretable, nor comparable across
studies, unless the treatment of the fault-onset window is stated explicitly.
A single frozen model differs by 0.085 in macro-F1 and 0.123 in MCC between two
defensible labeling conventions, which is larger than most of the improvements
reported between competing architectures on this benchmark. The effect is
moreover common to every model family tested, with a relative accuracy gain
between 0.142 and 0.150 across seven classifiers, so it displaces published
results without reordering them only when all studies happen to adopt the same
convention. We therefore
recommend that the treatment of the onset window be stated as part of the
experimental protocol in any observation-level study using these data.

Future work should investigate temporal aggregation, hierarchical diagnosis,
probability calibration, detection delay, and cost-sensitive alarm rules under
a new untouched test protocol. External validation with physical industrial
data will be necessary before deployment-oriented conclusions can be drawn.

\section*{Data and Code Availability}

The expanded Tennessee Eastman simulation data are publicly available from
Harvard Dataverse \cite{rieth2017data}. The CSV conversion used in this study is
available through Kaggle \cite{kaggleTEP}. All code, configuration files,
run-level split manifests, frozen hyperparameters, five-seed validation
outputs, final predictions, class-level metrics, confusion matrices, SHA-256
data hashes, and the environment specification are deposited at
an anonymized public repository for peer review, \TBD{anonymous Zenodo DOI};
the non-anonymized repository will be cited in the accepted version. The analysis of the
fault-onset window can be reproduced from the deposited predictions without
refitting any model.

% Acknowledgments removed for double-anonymized review.
% Restore in the accepted version.

\section*{Declaration of Competing Interest}

The authors declare that they have no known competing financial interests or
personal relationships that could have appeared to influence the work reported
in this paper.

\section*{CRediT Authorship Contribution Statement}

\textbf{Author 1:} Conceptualization, Methodology, Software, Validation,
Formal analysis, Investigation, Data curation, Writing -- original draft,
Visualization.
\textbf{Author 2:} Methodology, Software, Validation, Investigation,
Writing -- review \& editing.
\textbf{Author 3:} Conceptualization, Methodology, Resources, Writing --
review \& editing, Supervision, Project administration.

\begin{thebibliography}{99}

\bibitem{venkatasubramanian2003review}
V. Venkatasubramanian, R. Rengaswamy, K. Yin, and S. N. Kavuri,
``A review of process fault detection and diagnosis: Part I: Quantitative
model-based methods,'' \emph{Computers \& Chemical Engineering}, vol. 27,
no. 3, pp. 293--311, 2003. doi: 10.1016/S0098-1354(02)00160-6.

\bibitem{yin2014review}
S. Yin, S. X. Ding, X. Xie, and H. Luo,
``A review on basic data-driven approaches for industrial process monitoring,''
\emph{IEEE Transactions on Industrial Electronics}, vol. 61, no. 11,
pp. 6418--6428, 2014. doi: 10.1109/TIE.2014.2301773.

\bibitem{ji2022review}
C. Ji and W. Sun,
``A review on data-driven process monitoring methods: Characterization and
mining of industrial data,'' \emph{Processes}, vol. 10, no. 2, art. 335,
2022. doi: 10.3390/pr10020335.

\bibitem{chiang2001book}
L. H. Chiang, E. L. Russell, and R. D. Braatz,
\emph{Fault Detection and Diagnosis in Industrial Systems}. London, UK:
Springer, 2001. doi: 10.1007/978-1-4471-0347-9.

\bibitem{chiang2000discriminant}
L. H. Chiang, E. L. Russell, and R. D. Braatz,
``Fault diagnosis in chemical processes using Fisher discriminant analysis,
discriminant partial least squares, and principal component analysis,''
\emph{Chemometrics and Intelligent Laboratory Systems}, vol. 50, no. 2,
pp. 243--252, 2000. doi: 10.1016/S0169-7439(99)00061-1.

\bibitem{downs1993tep}
J. J. Downs and E. F. Vogel,
``A plant-wide industrial process control problem,''
\emph{Computers \& Chemical Engineering}, vol. 17, no. 3, pp. 245--255,
1993. doi: 10.1016/0098-1354(93)80018-I.

\bibitem{rieth2017data}
C. A. Rieth, B. D. Amsel, R. Tran, and M. B. Cook,
``Additional Tennessee Eastman Process simulation data for anomaly detection
evaluation,'' Harvard Dataverse, V1, 2017.
doi: 10.7910/DVN/6C3JR1.

\bibitem{breiman1984cart}
L. Breiman, J. H. Friedman, R. A. Olshen, and C. J. Stone,
\emph{Classification and Regression Trees}. Belmont, CA, USA: Wadsworth, 1984.

\bibitem{breiman2001rf}
L. Breiman, ``Random forests,'' \emph{Machine Learning}, vol. 45,
pp. 5--32, 2001. doi: 10.1023/A:1010933404324.

\bibitem{friedman2001gbm}
J. H. Friedman, ``Greedy function approximation: A gradient boosting machine,''
\emph{The Annals of Statistics}, vol. 29, no. 5, pp. 1189--1232, 2001.
doi: 10.1214/aos/1013203451.

\bibitem{chen2016xgboost}
T. Chen and C. Guestrin,
``XGBoost: A scalable tree boosting system,'' in \emph{Proceedings of the 22nd
ACM SIGKDD International Conference on Knowledge Discovery and Data Mining},
2016, pp. 785--794. doi: 10.1145/2939672.2939785.

\bibitem{cortes1995svm}
C. Cortes and V. Vapnik, ``Support-vector networks,'' \emph{Machine Learning},
vol. 20, pp. 273--297, 1995. doi: 10.1007/BF00994018.

\bibitem{pedregosa2011sklearn}
F. Pedregosa et al., ``Scikit-learn: Machine learning in Python,''
\emph{Journal of Machine Learning Research}, vol. 12, pp. 2825--2830, 2011.

\bibitem{matthews1975mcc}
B. W. Matthews, ``Comparison of the predicted and observed secondary structure
of T4 phage lysozyme,'' \emph{Biochimica et Biophysica Acta}, vol. 405, no. 2,
pp. 442--451, 1975. doi: 10.1016/0005-2795(75)90109-9.

\bibitem{kaggleTEP}
A. Melo,
``Tennessee Eastman Process CSV: Process simulation data for anomaly detection
evaluation,'' Kaggle Dataset, accessed Aug. 2026. [Online]. Available:
\url{https://www.kaggle.com/datasets/afrniomelo/tep-csv}.

% ---- NEW REFERENCES ADDED IN REVISION -------------------------
% VERIFY volume/page/DOI of each entry below against the publisher
% record before submission.

\bibitem{rieth2018issues}
C. A. Rieth, B. D. Amsel, R. Tran, and M. B. Cook,
``Issues and advances in anomaly detection evaluation for joint
human-automated systems,'' in \emph{Advances in Human Factors in Robots and
Unmanned Systems} (AHFE 2017), Cham: Springer, 2018, pp. 52--63.

\bibitem{reinartz2021extended}
C. Reinartz, M. Kulahci, and O. Ravn,
``An extended Tennessee Eastman simulation dataset for fault-detection and
decision support systems,'' \emph{Computers \& Chemical Engineering}, vol. 149,
art. 107281, 2021. doi: 10.1016/j.compchemeng.2021.107281.

\bibitem{yin2012comparison}
S. Yin, S. X. Ding, A. Haghani, H. Hao, and P. Zhang,
``A comparison study of basic data-driven fault diagnosis and process
monitoring methods on the benchmark Tennessee Eastman process,''
\emph{Journal of Process Control}, vol. 22, no. 9, pp. 1567--1581, 2012.
doi: 10.1016/j.jprocont.2012.06.009.

\bibitem{demsar2006statistical}
J. Dem\v{s}ar, ``Statistical comparisons of classifiers over multiple data
sets,'' \emph{Journal of Machine Learning Research}, vol. 7, pp. 1--30, 2006.

\end{thebibliography}

\end{document}
